\documentclass[11pt]{article}
\usepackage[margin=1in]{geometry}
\usepackage{amsmath,amssymb,amsthm}
\usepackage{graphicx}
\usepackage[hidelinks]{hyperref}

\newtheorem{theorem}{Theorem}
\newtheorem{lemma}{Lemma}
\newtheorem{remark}{Remark}

\newcommand{\C}{\mathbb C}
\newcommand{\R}{\mathbb R}
\newcommand{\ellTwo}{\ell_2}
\newcommand{\Sinf}{S_\infty}
\newcommand{\norm}[1]{\left\lVert #1\right\rVert}

\title{A Partial Extension of Ray's \(B_{p,q}\) Non-Embedding Criterion}
\author{FA/Banach discovery run \texttt{fa\_banach\_001}}
\date{June 13, 2026}

\begin{document}
\maketitle

\section{Source and Target}

The source paper is Samya Kumar Ray, ``On Isometric Embedding
\(\ell_p^m\to S_\infty^n\) and Unique operator space structure'',
arXiv:1911.00241.  For \(1\leq p,q<\infty\), Ray writes
\[
  B_{p,q}=\{(z_1,z_2)\in\C^2: |z_1|^p+|z_2|^q<1\}
\]
and denotes by \(\|\cdot\|_{B_{p,q}}\) the norm whose open unit ball is
\(B_{p,q}\).  Remark 4.2 proves non-embedding into \(\Sinf\) for
\(2<p,q<\infty\), records the positive triangular embedding for
\(B_{1,2}\), and states the conjecture that
\((\C^2,\|\cdot\|_{B_{p,q}})\) embeds isometrically into \(\Sinf\) if and
only if
\[
  (p,q)=(2,2),\quad (p,q)=(1,2),\quad\text{or}\quad (p,q)=(2,1).
\]

\begin{center}
\includegraphics[width=\textwidth]{figures/open_problem_crop.png}
\end{center}

\section{Candidate Partial Result}

\begin{theorem}
Let \(1\leq p,q<\infty\).  If \(\max\{p,q\}>2\), then
\((\C^2,\|\cdot\|_{B_{p,q}})\) does not embed linearly isometrically into
\(\Sinf\).
\end{theorem}

This is a partial result toward Ray's conjecture.  It extends the printed
\(2<p,q<\infty\) consequence in Remark 4.2 to every pair with at least one
exponent greater than \(2\).

\section{Proof}

We use the general criterion implicit in Ray's proof of Theorem 1.2 and
made explicit in Remark 4.2.  Let \(X=(\C^2,\|\cdot\|)\).  Suppose:
\begin{enumerate}
\item \(\|(z_1,z_2)\|=\|(\overline z_1,\overline z_2)\|\);
\item for every \(0<c\leq 1\), the map
  \((z_1,z_2)\mapsto (z_1,cz_2)\) is not a contraction from \(X\) to
  \(\ell_2^2\);
\item the set \(\{(z_1,z_2):\|(z_1,z_2)\|>|z_1|\}\) is dense in \(\C^2\);
\item \(\|(1,0)\|=\|(0,1)\|=1\);
\item \(\|(1,z)\|\geq 1\) and \(\|(w,1)\|\geq 1\) for all \(z,w\in\C\).
\end{enumerate}
Then Ray's descent argument shows that \(X\) cannot embed isometrically into
\(\Sinf\).  For completeness, we recall the mechanism.  If
\(\Phi(z_1,z_2)=z_1T+z_2S\) were an isometry into compact operators, then,
after the standard adjoint doubling and unitary conjugacy from Ray's proof,
one may assume that \(T\) is positive diagonal and has a finite-dimensional
one-eigenspace.  Birkhoff--James orthogonality gives
\(\zeta\in M_T\) with \(\langle T\zeta,S\zeta\rangle=0\).  Put
\(c=\|S\zeta\|_2\).  Then
\[
  \|(z_1,z_2)\|\geq
  \|z_1T\zeta+z_2S\zeta\|_2
  =(|z_1|^2+c^2|z_2|^2)^{1/2}.
\]
If \(c>0\), this says that \((z_1,z_2)\mapsto(z_1,cz_2)\) is a contraction
from \(X\) to \(\ell_2^2\), contradicting item 2.  If \(c=0\), the row and
column deletion step in Ray's proof applies, using items 1 and 3, and
produces a new isometric embedding with the multiplicity of the eigenvalue
\(1\) of \(T\) reduced by one.  After finitely many iterations \(T\) has norm
strictly less than \(1\), contradicting item 4.  Thus no such embedding
exists.

It remains to check the criterion for \(B_{p,q}\) when \(q>2\).  Items 1, 4,
and 5 are immediate from the Reinhardt unit ball.  For item 3, if \(z_2\neq0\)
then \(\|(z_1,z_2)\|_{B_{p,q}}>|z_1|\): otherwise, after scaling to the unit
ball boundary, the first term already contributes at least \(1\) and the
second contributes a positive amount.  Hence the required set contains
\(\{z_2\neq0\}\), which is dense.

For item 2, fix \(0<c\leq1\).  Choose \(y>0\) and set
\[
  x=(1-y^q)^{1/p}.
\]
Then \((x,y)\) lies on the \(B_{p,q}\)-unit sphere.  As \(y\downarrow0\),
\[
  x^2+c^2y^2
  =(1-y^q)^{2/p}+c^2y^2
  =1-\frac{2}{p}y^q+o(y^q)+c^2y^2.
\]
Since \(q>2\), the positive term \(c^2y^2\) dominates the order \(y^q\)
loss for all sufficiently small \(y\).  Thus \(x^2+c^2y^2>1\), so
\((z_1,z_2)\mapsto(z_1,cz_2)\) is not contractive.

Therefore no isometric embedding exists when \(q>2\).  If \(p>2\), apply the
same argument after swapping the two coordinates, which identifies
\((\C^2,\|\cdot\|_{B_{p,q}})\) isometrically with
\((\C^2,\|\cdot\|_{B_{q,p}})\).  This proves the theorem.

\section{Verification Notes}

The script \texttt{code/check\_Bpq\_obstruction.py} performs two sanity
checks.  First, for representative values with \(q>2\), it finds explicit
small \(y\) witnesses to the non-contraction condition above.  Second, it
checks the known triangular positive case \(B_{1,2}\) using the closed-form
operator norm of \(\begin{pmatrix} b&a\\0&b\end{pmatrix}\).  These numerical
checks are not part of the proof; they only guard against a sign or exponent
mistake.

\section{Limitations and Novelty Status}

This packet does not settle Ray's full conjecture.  In particular, it does
not address the remaining region \(1\leq p,q\leq2\) except for cases already
handled in the source paper.  The novelty check is local to this run:
the registry, solution folders, and attempt notes contained no prior entry
for arXiv:1911.00241 or the \(B_{p,q}\) conjecture.  No broad external
literature search was performed for this partial result.

Human review should focus on whether Ray's Remark 4.2 is being used exactly
as its proof supports, and on the elementary \(q>2\) non-contraction
calculation.

\end{document}
